% =====================================================================
%  CHAPTER 3: 各类食物的营养价值
% =====================================================================
\chapter{各类食物的营养价值}
\setcounter{chapter}{3}
\chapterintro{
掌握各类食物的营养特点（谷类、大豆、蔬果、畜禽水产品、蛋类、乳类、坚果）。熟悉食物营养价值的定义及影响因素。了解食物成分数据库。
}

% ----- 第一节 食物营养价值的定义及影响因素 -----
\section{食物营养价值的定义及影响因素}

\begin{defbox}
\textbf{食物营养价值（Nutritional Value of Food）：}食物中所含营养素和能量满足人体营养需要的程度。营养价值的高低取决于食物中营养素种类是否齐全、数量是否充足、相互比例是否适宜，以及是否易于被消化、吸收和利用。
\end{defbox}

\subsection{食物营养价值评价的维度}

\begin{enumerate}[leftmargin=*]
    \item \textbf{营养素种类与含量：}营养素种类齐全、含量丰富者营养价值较高。
    \item \textbf{营养素质量：}主要看该营养素的构成模式与人体的接近程度（如蛋白质氨基酸模式）。
    \item \textbf{消化吸收利用率：}即使含量丰富，若消化吸收差，实际营养价值也有限。
    \item \textbf{功能成分/生物活性物质：}不仅满足基本需求，更关注维护健康、预防疾病的额外价值。
\end{enumerate}

\subsection{影响食物营养价值的因素}

\begin{keybox}
\textbf{影响食物营养素组成和含量的主要因素：}
\begin{enumerate}[leftmargin=*]
    \item \textbf{品种、产地和季节：}不同品种的同种食物营养素差异大（如不同品种大米蛋白质含量差3\%\~{}4\%）；地理环境和气候条件影响产地；收获季节决定成熟度。
    \item \textbf{栽培和饲养条件：}土壤肥沃程度、施肥种类、饲料质量均显著影响农产品营养素水平。
    \item \textbf{收获成熟度：}未充分成熟的一般营养价值较低，但过度成熟可能导致营养素流失或口感变差。
    \item \textbf{加工和烹调：}加工精度（如碾米精度对B族维生素的影响极大）、热加工的温度和时间、加工方式（水煮损失水溶性维生素，油炸增加脂肪并损失热不稳定维生素）。
    \item \textbf{储藏条件：}温度、湿度、时间、光照、氧气影响营养素（尤其是维生素C、维生素E、$\beta$-胡萝卜素等抗氧化性维生素）的保留。
\end{enumerate}
\end{keybox}

\begin{tipbox}
\textbf{食物成分数据库：}我国《中国食物成分表》由营养与食品安全所主编，是国民营养与健康状况评价、膳食指导、食物营养标签等的基础性工具。包括原料型食物和加工型食物两大板块。
\end{tipbox}

% ----- 第二节 谷类、薯类及杂豆 -----
\section{谷类、薯类及杂豆}

\begin{keybox}
\textbf{谷类（Cereals）：}我国居民膳食的\keyword{主食}，提供约50\%\~{}70\%的能量和蛋白质。主要包括稻米（大米）、小麦（面粉）、玉米、小米、燕麦、荞麦等。
\end{keybox}

\subsection{谷类的主要营养特点}

\begin{enumerate}[leftmargin=*]
    \item \textbf{碳水化合物：}含量最高，主要为淀粉（70\%\~{}80\%），是人体最经济、最主要的功能来源。淀粉分为直链淀粉和支链淀粉。含少量膳食纤维（全谷中较丰富）。
    \item \textbf{蛋白质：}含量8\%\~{}12\%，\textbf{谷类蛋白的第一限制氨基酸为\keyword{赖氨酸}}。大米蛋白质含量最低（7\%\~{}8\%），生物价值在谷类中相对最高——因其赖氨酸含量稍优于其他谷类。
    \item \textbf{脂类：}含量低（1\%\~{}4\%），主要集中在胚芽（Germ）和糊粉层（Aleurone Layer）。谷类脂类主要含不饱和脂肪酸，且含有维生素E——这也是全谷物粉贮存中酸败的主要原因。
    \item \textbf{维生素：}谷类是B族维生素（尤其是硫胺素B1、尼克酸、吡哆醇B6）的重要来源，主要分布在糊粉层和胚芽。\textbf{碾磨精度越高，B族维生素损失越严重。}
    \item \textbf{矿物质：}约1.5\%\~{}3\%，以磷为主（多为植酸磷，吸收率低），钙、铁含量均较低。
\end{enumerate}

\subsection{全谷物与精制谷物的比较}

\begin{keybox}
\textbf{全谷物 vs 精制谷物：}
\begin{itemize}[leftmargin=*]
    \item \textbf{全谷物：}保留麸皮、胚乳和胚芽。富含膳食纤维、B族维生素、维生素E、矿物质和植物化学物（多酚、植物固醇、木酚素）。
    \item \textbf{精制谷物：}仅保留胚乳（淀粉+少量蛋白质），B族维生素、膳食纤维等大量损失。
    \item \textbf{推荐：}《中国居民膳食指南（2022）》建议每日摄入全谷物和杂豆50\~{}150 g。
\end{itemize}
\end{keybox}

\subsection{薯类}

\textbf{薯类：}马铃薯、甘薯、木薯等。主要营养特点：淀粉丰富（15\%\~{}25\%），蛋白质极低（1\%\~{}2\%），富含维生素C和钾，甘薯富含$\beta$-胡萝卜素（类胡萝卜素）。

\subsection{杂豆}

\textbf{杂豆（赤豆、绿豆、芸豆等）：}淀粉含量50\%\~{}60\%，蛋白质20\%\~{}25\%（高于谷类），膳食纤维丰富，B族维生素和矿物质含量较高。

% ----- 第三节 大豆类及其制品 -----
\section{大豆类及其制品}

\begin{keybox}
\textbf{大豆（Soybean）：}含蛋白质35\%\~{}40\%，是植物性食物中蛋白质含量最高者，且为\keyword{优质植物蛋白}。含脂肪15\%\~{}20\%，主要为不饱和脂肪酸（亚油酸占50\%以上）。还含大豆异黄酮、大豆皂苷、大豆磷脂、大豆低聚糖等有益健康成分。
\end{keybox}

\subsection{大豆及制品的蛋白质}

\begin{itemize}[leftmargin=*]
    \item 大豆蛋白质氨基酸模式较好，属于优质蛋白。
    \item \textbf{第一限制氨基酸为\keyword{蛋氨酸（Met）}}——这是与谷类蛋白质（第一限制氨基酸为赖氨酸）互补的基础。
    \item \textbf{谷类+豆类互补作用：}谷类的赖氨酸不足→由大豆的蛋氨酸缺陷结构得到互补，显著提高蛋白质生物价。
    \item 大豆制成的传统豆制品：豆腐（蛋白质6\%\~{}8\%）、豆浆（蛋白质2\%\~{}3\%）、腐竹（蛋白质高达44\%\~{}50\%，脂肪含量也相对较高）。
\end{itemize}

\subsection{豆制品分类}

\begin{enumerate}[leftmargin=*]
    \item \textbf{非发酵豆制品：}豆腐、豆浆、豆腐干、腐竹、豆皮等。
    \item \textbf{发酵豆制品：}酱油、豆豉、豆瓣酱、腐乳、纳豆等。发酵可提高B族维生素含量和蛋白质消化率，部分产生具有生物活性的肽类。
\end{enumerate}

\subsection{大豆中其他有益成分}

\begin{itemize}[leftmargin=*]
    \item \textbf{大豆异黄酮：}具有植物雌激素作用和抗氧化作用（见第2章）。
    \item \textbf{大豆皂苷：}降血脂、抗肿瘤、免疫调节作用。
    \item \textbf{大豆磷脂：}降血脂、保护肝脏、促进神经系统发育。
    \item \textbf{大豆低聚糖（水苏糖、棉子糖）：}促进双歧杆菌增殖（益生元作用）。
\end{itemize}

\begin{tipbox}
\textbf{豆类的抗营养因子及处理：}大豆中含胰蛋白酶抑制剂（引起胰腺肥大）、植物血凝素（凝集红细胞）、脂肪氧化酶（引起豆腥味）。这些因子均可通过充分加热（煮沸\~{}30min以上）彻底灭活。因此，豆浆必须煮沸后食用。
\end{tipbox}

% ----- 第四节 蔬菜、水果类 -----
\section{蔬菜、水果类}

\begin{keybox}
\textbf{蔬菜和水果的共同营养贡献：}蔬菜和水果是维生素C、类胡萝卜素、叶酸、维生素K、钾、镁和膳食纤维的最重要来源。同时富含多种植物化学物（多酚类、有机硫化物、类胡萝卜素等），对健康维护具有不可替代的重要意义。
\end{keybox}

\subsection{蔬菜的营养特点}

\begin{enumerate}[leftmargin=*]
    \item \textbf{碳水化合物：}含量因种类差异大（2\%\~{}20\%），主要为可溶性糖（果糖、葡萄糖、蔗糖）和不溶性膳食纤维（纤维素、半纤维素、木质素）。
    \item \textbf{蛋白质和脂肪：}含量极低（蛋白质≤2\%，脂肪<0.5\%）。
    \item \textbf{维生素：}
    \begin{itemize}[leftmargin=*]
        \item \keyword{维生素C}：深绿色叶菜、辣椒、西兰花、番茄中含量丰富。
        \item \keyword{类胡萝卜素}（维生素A原）：深绿色和橙黄色蔬菜中丰富（胡萝卜、菠菜、南瓜）。
        \item \keyword{叶酸（Folate）}：深绿色叶菜（菠菜、西兰花）、豆类和坚果中含量丰富。
        \item \keyword{维生素K}：绿叶蔬菜是最主要的膳食来源。
    \end{itemize}
    \item \textbf{矿物质：}富含钾（有助于维持血压正常）、镁。钙在部分蔬菜（如芥菜、苋菜）中含量较高，但草酸等抗营养因子影响其吸收。
    \item \textbf{植物化学物：}有机硫化物（葱属和十字花科）、多酚类（各种有色蔬菜）、皂苷等。
\end{enumerate}

\subsection{水果的营养特点}

\begin{enumerate}[leftmargin=*]
    \item \textbf{碳水化合物：}主要为简单糖（果糖、葡萄糖、蔗糖），含量6\%\~{}25\%。含水量一般80\%以上，能量较低。
    \item \textbf{维生素：}新鲜水果是\keyword{维生素C}的最主要来源（猕猴桃、柑橘、草莓、鲜枣中含量极高）。也含类胡萝卜素（芒果、杏、木瓜）。
    \item \textbf{矿物质：}钾含量丰富（与血压控制相关）。
    \item \textbf{膳食纤维：}果胶含量较高，可降低血清胆固醇。
    \item \textbf{植物化学物：}花色素（浆果）、白藜芦醇（葡萄）、类胡萝卜素、黄酮类等。
\end{enumerate}

\begin{tipbox}
\textbf{蔬菜与水果的摄入建议（膳食指南2022）：}蔬菜摄入量300\~{}500 g/d，其中深色蔬菜应占一半（\~{}250 g/d）；水果200\~{}350 g/d，不能以果汁替代整果。蔬菜和水果在营养素谱系上有互补性，不可相互替代。
\end{tipbox}

\begin{keybox}
\textbf{蔬菜营养损失的关键因素：}
\begin{itemize}[leftmargin=*]
    \item \textbf{水溶性维生素：}维生素C和B族维生素在洗切、浸泡和煮沸中极易损失。
    \item \textbf{热敏性维生素：}维生素C、叶酸和部分B族维生素在高温下大量破坏。
    \item \textbf{氧化损失：}切碎后接触空气中的氧和光照，加速维生素氧化破坏。
    \item \textbf{合理烹调建议：}先洗后切、急火快炒、炒完即食，减少加工损失。
\end{itemize}
\end{keybox}

% ----- 第五节 畜、禽、水产品 -----
\section{畜、禽、水产品}

\begin{keybox}
\textbf{动物性食物的核心营养贡献：}畜、禽和水产品是\keyword{优质蛋白质}、\keyword{血红素铁}和\keyword{维生素B${}_{12}$}的最重要来源。B12是唯一只能从动物性食物中获得的维生素。
\end{keybox}

\subsection{畜肉（猪、牛、羊）}

\begin{itemize}[leftmargin=*]
    \item \textbf{蛋白质：}含量15\%\~{}22\%，氨基酸模式与人体接近，为优质蛋白。
    \item \textbf{脂肪：}含量因种类、部位而异（5\%\~{}30\%），以\keyword{饱和脂肪酸}为主（约占40\%\~{}60\%）。内脏脂肪含量更高，且胆固醇含量高（猪脑、猪肝尤其突出）。
    \item \textbf{碳水化合物：}极低（<1\%，主要以糖原形式存在）。
    \item \textbf{矿物质：}\textbf{血红素铁（Heme Iron）}——吸收率高达15\%\~{}25\%，是植物性非血红素铁（吸收率1\%\~{}5\%）的数倍。锌的吸收利用率也高。
    \item \textbf{维生素：}B族维生素（硫胺素B1——猪肉含量最高；吡哆醇B6——鸡肉中丰富；B12——所有肉类中均有）。维生素A在内脏（肝）中极高。
\end{itemize}

\subsection{禽肉（鸡、鸭）}

\begin{itemize}[leftmargin=*]
    \item 蛋白质含量与畜肉相当，但\keyword{不饱和脂肪酸}比例高于畜肉、胆固醇含量较低。
    \item 鸡胸肉蛋白质含量最高、脂肪最少（\~{}20\%蛋白、3\%脂肪），为健身人群的首选动物蛋白来源。
\end{itemize}

\subsection{水产品（鱼、虾、贝类）}

\begin{itemize}[leftmargin=*]
    \item \textbf{蛋白质：}与畜肉相当，但更易消化。
    \item \textbf{脂肪：}显著低（大部分<5\%），以\keyword{不饱和脂肪酸}为主——尤其是\keyword{n-3多不饱和脂肪酸}$\Rightarrow$深海鱼油含EPA（二十碳五烯酸）和DHA（二十二碳六烯酸），对心血管保护和脑发育有重要作用。
    \item \textbf{矿物质：}海产品富含碘、锌、硒（牡蛎——"海中的锌宝库"）。
    \item \textbf{维生素：}鱼肝含丰富的维生素A和维生素D。
\end{itemize}

\begin{exambox}
\textbf{考点提示——动物铁 vs 植物铁：}
\begin{itemize}[leftmargin=*]
    \item 动物食物：血红素铁，直接进入小肠上皮细胞，吸收率15\%\~{}25\%，不受膳食中植酸等抑制因素影响。
    \item 植物食物：非血红素铁，需先转化为Fe$^{2+}$，吸收率1\%\~{}5\%，受植酸、多酚、钙等抑制因素影响极大。
\end{itemize}
这是营养学中判断食物铁营养价值的关键区别。
\end{exambox}

% ----- 第六节 蛋类及蛋制品 -----
\section{蛋类及蛋制品}

\begin{keybox}
\textbf{蛋类的营养地位：}全鸡蛋蛋白质的\keyword{氨基酸模式与人体最接近}——是自然食物中生物价值最高的蛋白质来源。常被作为"参考蛋白质（Reference Protein）"用于食物蛋白质质量评价。
\end{keybox}

\subsection{蛋类的化学成分}

\textbf{全蛋：}
\begin{itemize}[leftmargin=*]
    \item 蛋白质11\%\~{}14\%，大部分为卵清蛋白（蛋清中），8种必需氨基酸齐全，生物价（BV）高达94。
    \item 脂肪8\%\~{}11\%，集中于蛋黄中，主要为甘油三酯和磷脂（卵磷脂为主），不饱和脂肪酸为主。
    \item 胆固醇含量高：蛋黄中约1500 mg/100g（每个鸡蛋约200\~{}300 mg胆固醇）。
    \item 微量营养素：维生素A、D、B2、B12、生物素含量较高；矿物质以磷、钙、铁、锌和硒为主（蛋黄中）。
\end{itemize}

\textbf{蛋清 vs 蛋黄：}
\begin{table}[h!]
\centering
\begin{tabular}{p{2.5cm}p{4cm}p{4cm}}
\hline
\textbf{营养成分} & \textbf{蛋清（约占蛋重60\%）} & \textbf{蛋黄（约占蛋重30\%）} \\
\hline
蛋白质 & 以卵清蛋白为主，约10\% & 卵黄蛋白为主，约16\% \\
脂肪 & 含量极低 & 30\%\~{}33\%，含丰富的卵磷脂 \\
维生素 & 极少 & 富含A、D、E、B族 \\
矿物质 & 极少 & 钙、磷、铁、锌、硒 \\
胆固醇 & 无 & 约1500 mg/100g \\
\hline
\end{tabular}
\end{table}

\subsection{蛋制品的营养变化}

\begin{itemize}[leftmargin=*]
    \item \textbf{咸蛋：}水分减少，钠含量显著增高，其他营养素基本保留。
    \item \textbf{皮蛋：}碱性加工条件下维生素B族大量破坏，维生素A和D可保留。
    \item \textbf{蛋粉：}脱水浓缩，营养素浓缩，但维生素A和D有一定损失。
\end{itemize}

% ----- 第七节 乳及乳制品 -----
\section{乳及乳制品}

\begin{keybox}
\textbf{乳类的核心营养地位：}乳类（特别是牛乳）是\keyword{最佳膳食钙来源}——不仅钙含量高（\~{}104 mg/100g），而且钙的吸收率高达32.1\%，远超其他钙源。乳类还是维生素B2（核黄素）和优质蛋白质的重要来源。
\end{keybox}

\subsection{牛乳的营养成分}

\begin{enumerate}[leftmargin=*]
    \item \textbf{蛋白质：}3.0\%\~{}3.5\%，主要为\keyword{酪蛋白（Casein）}（占80\%）和\keyword{乳清蛋白（Whey）}（占20\%）。乳清蛋白中的$\alpha$-乳白蛋白是母乳中的优势蛋白（见第4章）。生物价87，仅次于全鸡蛋。
    \item \textbf{脂肪：}3\%\~{}4\%，以甘油三酯为主，饱和脂肪酸为主，胆固醇含量低（10\~{}15 mg/100g）。以乳化小颗粒状存在，易于消化。
    \item \textbf{碳水化合物：}主要为\keyword{乳糖（Lactose）}（4.5\%\~{}5.0\%）。
    \begin{itemize}[leftmargin=*]
        \item 乳糖在肠道被乳糖酶水解为葡萄糖+半乳糖。
        \item \textbf{乳糖不耐症（Lactose Intolerance）：}因乳糖酶活性不足，未消化的乳糖进入结肠被细菌发酵产生气体和酸性产物，出现腹胀、腹痛和腹泻。亚洲人群发生率极高（80\%\~{}90\%）。
    \end{itemize}
    \item \textbf{钙：}含量104 mg/100g，吸收率32.1\%——高含量+高吸收率使牛乳成为不可替代的膳食钙源。
    \item \textbf{维生素：}维生素A（全脂乳丰富）、维生素B2（约0.13 mg/100g，是重要的核黄素来源）、维生素D（通常强化）。维生素C含量极低，在加热中大部分损失。
    \item \textbf{矿物质：}钙磷比为1.3:1，利于钙吸收。铁含量极低（<0.05 mg/100g，因此牛乳是缺铁性贫血的危险因素——婴幼儿单纯依靠牛乳可引致缺铁）。
\end{enumerate}

\subsection{常见乳制品}

\begin{itemize}[leftmargin=*]
    \item \textbf{酸奶（Yogurt）：}发酵后乳糖含量降低，可部分缓解乳糖不耐受；益生菌有利于肠道健康。
    \item \textbf{奶酪（Cheese）：}蛋白质和钙浓缩（钙400\~{}800 mg/100g），脂肪浓缩（25\%\~{}35\%），乳糖极低（适合乳糖不耐受者）。
    \item \textbf{奶粉（Milk Powder）：}脱水制成，营养素浓缩，适宜储藏和运输。
\end{itemize}

% ----- 第八节 坚果类 -----
\section{坚果类}

\begin{keybox}
\textbf{坚果的营养特点：}坚果富含不饱和脂肪酸、优质蛋白质、维生素E、B族维生素和多种矿物质（镁、锌、硒），是能量和营养素的浓缩食物。主要分为两类：树坚果（核桃、杏仁、腰果等）和籽类（葵花籽、南瓜籽等）。
\end{keybox}

\subsection{坚果的主要营养价值}

\begin{enumerate}[leftmargin=*]
    \item \textbf{脂肪：}含量极高（\~{}50\%，有的高达60\%+），以不饱和脂肪酸（单不饱和+多不饱和）为主。核桃的$\alpha$-亚麻酸（n-3必需脂肪酸）含量在食物中名列前茅。
    \item \textbf{蛋白质：}15\%\~{}25\%，大豆除外，是植物性食物中的高蛋白来源。
    \item \textbf{维生素：}维生素E含量丰富（重要的脂溶性抗氧化剂）、B族维生素（尤其是叶酸和吡哆醇B6）。
    \item \textbf{矿物质：}镁、钾、铜、锌、硒含量较高。
    \item \textbf{植物化学物：}植物固醇、多酚类、木酚素等。
\end{enumerate}

\subsection{食用注意事项}

\begin{itemize}[leftmargin=*]
    \item \textbf{能量密度极高：}100 g核桃约600\~{}700 kcal，过量摄入易导致能量过剩和体重增加。
    \item \textbf{中国居民膳食指南（2022）：}每日建议摄入量10\~{}15 g（约一小把），不宜贪多。
    \item \textbf{过敏风险：}花生和坚果是最常见的食物过敏原之一，在婴幼儿添加食物中须注意。
\end{itemize}

% ----- 复习题 -----
\reviewcontent{
\section{复习题}

\begin{enumerate}[leftmargin=*, itemsep=8pt]
    \item \textbf{【名词解释】}食物营养价值、优质蛋白质、血红素铁、乳糖不耐症、氨基酸模式、限制氨基酸、蛋白质互补作用
    \item \textbf{【简答题】}影响食物营养价值的因素有哪些？
    \item \textbf{【简答题】}谷类食物的主要营养特点是什么？为什么精制谷物不如全谷物？
    \item \textbf{【简答题】}简述大豆蛋白质的营养特点，及其与谷类蛋白质互补作用的原理。
    \item \textbf{【简答题】}试述畜肉、禽肉和水产品在蛋白质、脂肪和铁含量上的主要差异。
    \item \textbf{【简答题】}为什么全鸡蛋蛋白质被称为"参考蛋白质"？鸡蛋中的微量营养素集中在哪部分？
    \item \textbf{【简答题】}为什么牛乳被认为是最佳膳食钙来源？乳糖不耐受症的机制是什么？
    \item \textbf{【简答题】}简述蔬菜和水果在维生素C、类胡萝卜素和膳食纤维上的共同营养贡献，以及二者的差异。
    \item \textbf{【论述题】}从蛋白质、脂肪、矿物质和生物活性成分四个维度，比较分析畜肉类、水产品、蛋类和奶类的营养价值差异。
\end{enumerate}

\subsection*{参考答案要点}

\begin{enumerate}[leftmargin=*, itemsep=5pt]
    \item 见本章各概念定义框。
    \item 品种/产地/季节、栽培和饲养条件、收获成熟度、加工/烹调、储藏条件。
    \item 谷类：碳水化合物70\%\~{}80\%（最经济的热能来源），蛋白质8\%\~{}12\%（赖氨酸为第一限制氨基酸），低脂，B族维生素（精加工损失严重）。全谷物优于精制：多含膳食纤维、B族维生素、维生素E、矿物质和植物化学物。
    \item 大豆蛋白35\%\~{}40\%，优质植物蛋白，第一限制氨基酸为蛋氨酸；谷类第一限制氨基酸为赖氨酸。二者混合，谷类提供蛋氨酸、豆类提供赖氨酸，实现必需氨基酸互补。
    \item 蛋白质均优质。脂肪：畜肉饱和脂肪酸多，禽肉不饱和脂肪酸比例较高，水产品以n-3多不饱和脂肪酸为特色。铁：三者均含血红素铁，吸收率15\%\~{}25\%。
    \item 全鸡蛋氨基酸模式与人体最接近，BV高达94。微量营养素集中在蛋黄中（维生素A、D、E、B12、磷、铁、锌、硒、胆固醇）。
    \item 钙含量高（104 mg/100g），吸收率高（32.1\%）。乳糖不耐受机制：乳糖酶活性不足，未被消化的乳糖在结肠经细菌发酵产酸产气。
    \item 共同贡献：维生素C、类胡萝卜素、叶酸、钾、镁、膳食纤维、植物化学物。差异：蔬菜密度更低、可作为主菜食用量大，水果含糖量（简单糖）高，不宜替代蔬菜。
    \item (1) 蛋白质：均优质，蛋类BV最高——参考蛋白质；(2) 脂肪：水产品最低，n-3不饱和脂肪酸具有心血管保护作用；(3) 矿物质：畜肉等动物肉含高吸收率的血红素铁和锌，乳类的钙最高吸收率；(4) 生物活性成分：水产品含EPA/DHA，乳含益生菌/免疫球蛋白，蛋含卵磷脂。
\end{enumerate}

\newpage
}
